\clearpage
\heading{K.4. Moving the named constants outside a run}\label{sec:15k4}
\smallref{Pass 15A. \paper{Lemma K.4, p.~57}. \pub{NS1/Machine.lean}{205}{Run data}, \pub{NS1/Machine.lean}{229}{evaluation rules}, \pub{Paper/Postpone.lean}{30}{quantified statement}.}

Let $A,R$ be commutative rings with structure map $\iota:A\to R$.
Fix the base-weight function $\theta$, a program $p$, its argument $z$, and a
derivation that every run of $p(z)$ halts. It determines an ordered finite
list $L$ of accepting occurrences. Each $r\in L$ has cost $c_r$, base
weight $a_r\in A$, an ordered bill $b_r$ of binary names, and an output
tree. With a name valuation $\kappa$ the weighted sum is
\[
 V=\sum_{r\in L}\iota(a_r)\prod_{s\in b_r}\kappa(s).
\]
A halted rejection contributes no occurrence. A zero-weight acceptance
still contributes an occurrence. Total halting is a separate premise. The Lean function returning the accepting list defaults to the empty list
when a halting derivation does not exist. An empty list by itself cannot
witness that a computation halted.

Assume $c_r\le T<2^J$ for every accepting occurrence. A \emph{use} adds
one name and costs one extra step; sequential composition concatenates
the bills and adds both costs. Induction over the evaluation rules gives
$|b_r|\le c_r$. Index the leaves of a trie with $Q'$ nodes and height
$\eta$ by names $s_1,\ldots,s_m$, and write $\lambda_i$ for their
values. Under the manuscript's leaf-membership hypothesis, the numbers
$n_i(r)=\#\{s_i\text{ in }b_r\}$ all fit in $J$ bits.

\subhead{What one tally accepts}
For $e\in\{0,\ldots,2^J-1\}^m$, let $z_e=(\mathrm{WF},(p,(z,e)))$,
with each leaf counter stored least significant bit first. The interpreter
follows the original control flow, including calls to computed programs.
Each \emph{use} becomes a deterministic decrement at its named leaf.
An absent leaf or an attempted decrement from zero rejects the branch.
Each successful decrement preserves the counter width and trie shape.
The terminal test accepts exactly when all residual counters are zero.
Thus the accepted occurrences for $e$ have $n_i(r)=e_i$ for all $i$.
Reading a base weight through $\theta$ still multiplies by the same
$a$-factor and issues no name.

The structural induction couples each surviving occurrence with the original
output and residual budget. The coupling respects concatenation of bills
and the order of guesses. The terminal test keeps one occurrence or zero.
Consequently the ordered list of base weights in $z_e$ equals
$[a_r\mid r\in L,\ n(r)=e]$, including duplicates. Every surviving
interpreter bill is empty. Under the leaf-membership hypothesis each
occurrence belongs to exactly one such list. Commutativity, distributivity,
and transport by $\iota$ now give
\[
 V=\sum_e \Bigl(\prod_i\lambda_i^{e_i}\Bigr)\,\iota\bigl(\mathrm{pv}(z_e)\bigr).
\]
\smallref{\pub{NS1/Accel.lean}{79}{Bill length}, \pub{NS1/Eval.lean}{213}{composition}, \pub{NS1/EvalThm.lean}{138}{interpreter induction}, \pub{NS1/Postpone.lean}{321}{ordered filtering}, \pub{NS1/Postpone.lean}{415}{finite sum identity}.}

\clearpage
\heading{Paying for the tally interpreter}
\smallref{\pub{NS1/Budget.lean}{68}{Budget height and size}, \pub{NS1/Interp.lean}{109}{decrement}, \pub{NS1/Interp.lean}{228}{terminal test}, \pub{NS1/Postpone.lean}{344}{complete interpreter cost}.}

Write $d_e$ for the height of the \emph{counter} trie and $n_e$ for its
number of nodes and digits. A root-to-leaf descent pays for at most
$\eta$ path bits plus $J$ counter digits and the leaf marker. Adding $J$
digits to each of $m$ leaves gives
\[
 d_e\le\eta+J+1,\qquad n_e=Q'+mJ\le Q'(J+1).
\]
Even when subtraction borrows through every bit, the fixed-width
little-endian routine takes at most $20J$ steps on a successful call.
Path descent, rebuilding the traversed nodes, and subtraction together
cost at most $60d_e$ on success. Since each update preserves both $d_e$
and $n_e$, later uses cost no more than earlier ones at the same depth.

The interpreter induction charges
$M=1000+100d_e$ per original step. For example the name instruction
needs $60d_e+420\le M$, including its surrounding calls. Its induction
hypothesis already paid for evaluating the expression producing the name.
The branch for a nested \emph{run} carries the residual counters into
the called computation. Thus its subcosts add within the original run
cost. Recursive calls introduce no new multiplicative depth factor.
After the simulation, a deterministic traversal checks every counter for
zero at cost at most $40n_e$ on a surviving branch. Including the fixed
wrapper instructions, an occurrence corresponding to $r$ costs at most
$Mc_r+40n_e+40$. The all-branches halting derivation and this per-acceptance
bound give precisely
\[
   \mathrm{Plain}_{T'}(z_e),\qquad
   T'=\bigl(1000+100(\eta+J+1)\bigr)T+40Q'(J+1)+40.
\]
The interpreter code is fixed. Three constructor nodes join it to $p$,
$z$, and the encoding of $e$. That encoding has at most $4n_e$ nodes.
Therefore
\[
 |z_e|\le |\mathrm{WF}|+|p|+|z|+4Q'(J+1)+3.
\]
This measures the description of one query. The evaluation theorem
gives a bound per accepting run and halting on every run. The full sum
may have $2^{Jm}$ tally vectors. Neither deterministic enumeration time
nor peak interpreter workspace follows from the displayed bounds.

\subhead{Where the two consumers need the bounds}
K.5 uses the description bound to license cutting on every short input.
On a plain outer input, K.4 supplies the plain postponed query needed for
the cutter's value clause. K.6 first bounds $|z_e|$ for each possible
self-query using the counter encoding alone. This allows execution of
the candidate evaluator on $z_e$ and hence proves total halting and an
accepting-time bound for the self-program. Applying K.4 to that program
then proves all the queries plain. Only after verifying $T'\le v^H$ can
K.6 use the evaluator's conditional value equality on each query.

\smallref{\pub{NS1/EvalThm.lean}{252}{Nested calls}, \pub{NS1/EvalThm.lean}{352}{name instruction cost}, \pub{NS1/Guess.lean}{187}{counter bounds}, \pub{NS1/Scale.lean}{93}{K.5 value step}, \pub{Paper/NoAcc.lean}{181}{K.6 postponement}.}

\clearpage
\heading{Which occurrences survive a zero contribution?}
\smallref{\pub{NS1/Postpone.lean}{48}{Totalized lookup}, \pub{NS1/Postpone.lean}{108}{leaf matching}, \pub{NS1/Postpone.lean}{141}{node matching}, \pub{NS1/Postpone.lean}{246}{one-bill identity}.}

The formal hypothesis requires $\kappa(s)=\mathrm{lookup}(s)$ on every
used name. It permits an unsupported request when both sides vanish. A lookup outside the trie is defined to be zero. A run asking
for such a name is rejected for every $e$. Its original product is zero,
so the weighted sum still agrees. A genuine leaf with value zero gives
different run semantics. Its occurrences survive for the matching $e$,
even though their external coefficient vanishes when the count is positive.
This distinction locates precisely where the printed hypothesis enters the
one-tally-per-run reading.

Consider the following ordered finite run list over $\mathbb Z$ with $J=2$.
The leaves $0,10,11$ carry $2,-1,0$. The internal path $1$ has lookup
zero. Entries in parentheses are lists of names.
\begin{center}
\begin{tabular}{ccll}
Occurrence & Base weight & Bill & Matching $e$ \\
\hline
1 & $2$  & $(0,10)$ & $(1,1,0)$ \\
2 & $7$  & $()$     & $(0,0,0)$ \\
3 & $2$  & $(10,0)$ & $(1,1,0)$ \\
4 & $-3$ & $(0,10)$ & $(1,1,0)$ \\
5 & $5$  & $(11)$   & $(0,0,1)$ \\
6 & $0$  & $(10,0)$ & $(1,1,0)$ \\
7 & $11$ & $(1)$    & none
\end{tabular}
\end{center}
The tally $(1,1,0)$ preserves the ordered base weights $[2,2,-3,0]$. Its coefficient is $2(-1)$. Including the empty bill and zero-valued leaf
gives $7-2(2+2-3+0)+0(5)=5$, as direct weighting does. Keeping only the
distinct base weights would change this value. Requiring the interpreter's
final output tree or cost to equal the original would exceed the public
lemma's conclusion. It compares the ordered weights after filtering.

At the level of arbitrary bills, width has a sharp limit. Four uses of leaf $0$ need counter
value four. For $J=2$ no counter can hold it, and a base-weight-one run
would lose its contribution $2^4=16$. Conversely, two uses each of
$0$ and $10$ fit at the same width. The total-bill cutoff is a convenient
sufficient condition. Individual counts are the obstruction in this
finite identity.

The independent finite check compares binary subtraction against exact
count vectors for the above three leaves, allowing the empty path, an
internal path, and an overlong path as adversarial names. It checks
16,576 bill--budget pairs, including 40 supported bills and 219 bills
with an absent name. It also checks both four-use boundary examples,
width zero, and the ordered signed run list. These are finite checks of
the tally algebra. The arbitrary-program halting and machine-cost
statements use the pinned Lean theorems.
\smallref{Reproduce with \texttt{python3 checks/pass15\_k4.py} from the review root.}

\clearpage
\heading{The exact time endpoint for a fixed counter width}\label{sec:15endpoint}
The four-use example raises a question about the program that produced the bill. A mathematical list can have four entries. Can an accepting run of the actual machine issue four requests in four steps? Its first step already supplies a unit of slack.

\subhead{A structural invariant for accepting runs}
For every evaluation derivation and every accepting occurrence $r$, define the integer $d(r)=\mathrm{cost}(r)-|\mathrm{bill}(r)|$. Then $d(r)\ge1$. At a quote, an identity or either outcome of a guess the cost is one and the bill is empty. A failed computation has no accepting occurrence. A projection or base-weight instruction adds a step and keeps the same bill length. A named multiplication adds one step and one request, preserving $d$. A pair or conditional combines the computations taken along that branch. Its bills concatenate and it pays an additional step. A computed call combines the computations producing the code and argument with the execution of the callee, and also pays an additional step. These are all the evaluation constructors. Induction gives $|\mathrm{bill}(r)|<\mathrm{cost}(r)$ for each accept, including those of zero weight.

Consequently the inclusive clock $\mathrm{cost}(r)\le T\le 2^J$ gives $|\mathrm{bill}(r)|<2^J$. This is exactly the premise of the one-bill summation identity. The wrapper analysis used only $\mathrm{cost}(r)\le T$. The full K.4 conclusion therefore extends to $T=2^J$, with the same time bound, encoded size, scalar sum and ordered filter identity. The publisher uses the strict clock $T<2^J$. Its choices in K.5 and K.6 already satisfy that premise.

\subhead{One more step can overflow}
Build $F_0$ by quoting the empty tree and $F_{n+1}$ by applying one name instruction to $F_n$. On the empty argument its unique accepting record has cost $n+1$, base weight one, output the empty tree and $n$ copies of the empty word as its bill. The induction is immediate from the machine's name rule, which outputs the empty tree again at each step. Take $n=2^J$ and put the integer $2$ at the trie's single root leaf. The original value is $2^{2^J}$. No $J$-bit tally can match this bill, so every wrapper rejects and the expanded sum is zero. Thus $2^J$ is the largest uniform accepting clock for this width with unrestricted tries and programs. This example includes $J=0$.

\smallref{Publisher rules: \pub{NS1/Machine.lean}{230}{evaluation constructors}, \pub{NS1/Accel.lean}{79}{weak bill-length bound}, \pub{NS1/Postpone.lean}{415}{summation premise}. The strengthened bound, endpoint and family of runs are review lemmas in \texttt{checks/pass15\_plain.lean}.}
